\section{Related Work}
\label{sec:related}

\para{Plan-and-write story generation.}
The idea of separating narrative planning from prose generation dates to early neural storytelling systems.
\citet{yao2019planandwrite} introduced a two-stage framework that extracts keyword storylines via RAKE and conditions generation on these plans, achieving significant coherence improvements on \rocstories.
\citet{fan2019strategies} extended this to longer narratives using SRL-based action sequences and entity-anonymized intermediate representations.
PlotMachines~\cite{rashkin2020plotmachines} added dynamic state tracking over outline elements, while \citet{goldfarb2020aristotelian} introduced Aristotelian rescoring classifiers for principled plot quality control.
These works establish that decomposing generation into planning and writing stages improves coherence, but they rely on shallow structural representations (keywords, SRL predicates) rather than rich narrative components.

\para{Hierarchical and multi-agent decomposition.}
Recent work uses deeper decompositions with LLM-based planning.
DOC~\cite{yang2023doc} generates hierarchical outline trees with three levels of plot events, character state tracking, and FUDGE-based adherence control, achieving large coherence and relevance gains over prior systems.
Agents' Room~\cite{huot2024agentsroom} takes decomposition further by assigning specialized agents to distinct narrative dimensions (conflict, character, setting, plot) and writing stages (exposition through resolution), with a shared scratchpad for coordination.
DSR~\cite{dsr2025} decomposes screenwriting into narrative creation and format conversion, introducing ``Narrative Directives'' as chain-of-thought reasoning over pacing, exposition, and character arcs.
These systems consistently outperform end-to-end generation, but each uses ad-hoc decompositions tied to their specific pipelines.
Unlike these approaches, we treat story components as a \emph{composable vocabulary} and systematically evaluate the effect of combinatorial recombination.

\para{Synthetic data for creative writing.}
Several methods generate synthetic training data for narrative tasks.
Weaver~\cite{xie2024weaver} uses instruction backtranslation---collecting human-written text and using GPT-4 to generate corresponding instructions---to create 500K instruction-response pairs for creative writing.
DSR~\cite{dsr2025} employs a hybrid approach combining reverse compression of inputs with forward distillation of outputs, showing that reverse synthesis alone creates input-output misalignment.
Agents' Room~\cite{huot2024agentsroom} uses teacher-student backtranslation from Gemini Ultra to Gemini Flash.
Our work is closest in spirit to these approaches but uses a different mechanism: we extract structured components and recombine them, analogous to how \skillmix creates novel skill combinations for instruction-following data.

\para{Skill-based synthetic data generation.}
\citet{allenzhu2024instructskillmix} introduced the \skillmix paradigm for general instruction-following: enumerate skills, create combinatorial mixtures, generate synthetic training data, and fine-tune models on this data.
This approach significantly improved instruction-following capabilities across diverse tasks.
We test whether this paradigm transfers to the harder domain of creative writing, where ``skills'' correspond to narrative components and quality depends on subjective aesthetic judgments rather than objective correctness.

\para{LLM creative writing limitations.}
\citet{chakrabarty2023art} documented that LLM-generated stories are less diverse and more formulaic than human-written ones, with limited vocabulary variation and predictable narrative patterns.
CRITICS~\cite{bae2024critics} addresses this through multi-critic iterative refinement, using persona-assigned LLM critics to evaluate and improve creativity dimensions.
Re3~\cite{yang2022re3} uses recursive reprompting and revision for long story generation.
Our work connects to this literature by testing whether structural decomposition---an orthogonal strategy to iterative refinement---can mitigate the formulaic ceiling.
